\section{Algorithm: Semantic-Similarity Per-Turn Router}
\label{sec:routing}

\subsection{Approach}

The per-turn router assigns each main-thread turn to a tier in
\(\tierset = \{\text{haiku}, \text{sonnet}, \text{opus}\}\) without
access to role labels or structured signal.  We adopt a
\emph{semantic-similarity} approach: embed the turn into a shared
representation space, retrieve the most similar labelled exemplars,
and return the majority-vote tier of the top-\(k\) neighbours.

The exemplar set contains 50 turns per tier (150 total), drawn from
real \blackrim main-thread session transcripts and annotated following
the tier-boundary rubric in \texttt{docs/research/model-cost-quality-landscape.md}
(annotation: $\sim$2--4 hours).  Turns are encoded as the concatenation of
\texttt{user\_prompt} and \texttt{observed\_response\_summary} using
\texttt{all-MiniLM-L6-v2} (L2-normalised, \(d\!=\!384\); selected by RC-04
bake-off; swappable via \texttt{SEMANTIC\_ROUTER\_MODEL}).

At construction time, each exemplar is encoded once; the embedding
matrix \(\mathbf{E} \in \mathbb{R}^{|E| \times d}\) persists for the
session lifetime.

\subsection{Algorithm}

\begin{algorithm}[H]
\SetAlgoLined
\DontPrintSemicolon
\KwIn{turn string $s(\turn_i)$,
      exemplar matrix $\mathbf{E}$, exemplar tiers $\{t_e\}$,
      $k$, $\tau$, escalation tier $\tier_{\text{esc}} = \text{opus}$}
\KwOut{predicted tier $\tier_i$}
\BlankLine
$\mathbf{v} \gets \phi\bigl(s(\turn_i)\bigr)$
\tcp*[r]{embed query (L2-normalised)}
$\mathbf{s} \gets \mathbf{E}\,\mathbf{v}$
\tcp*[r]{cosine similarities, $\mathbf{s} \in \mathbb{R}^{|E|}$}
\If{$\max(\mathbf{s}) < \tau$}{
  \Return{$\tier_{\text{esc}}$}
  \tcp*[r]{conservative escalation: OOD turn}
}
$I \gets \operatorname{top-}k(\mathbf{s})$
\tcp*[r]{indices of $k$ largest similarities}
$\hat\tier \gets \operatorname{majority\_vote}\bigl(\{t_e : e \in I\}\bigr)$
\tcp*[r]{ties resolved toward higher tier}
\Return{$\hat\tier$}\;
\caption{Semantic-similarity per-turn router (conservative escalation)}
\label{alg:router}
\end{algorithm}

Default parameters: \(k = 5\), \(\tau = 0.30\).  Tie-breaking in
majority vote prefers the more capable tier (opus over sonnet over
haiku), making the router conservative by design.

\subsection{Inference overhead}

Encoding a 512-token turn with \texttt{all-MiniLM-L6-v2} on CPU takes
approximately 10\,ms; the matrix-vector product over 150 exemplars
adds under 1\,ms.  Total routing overhead is $\sim$10--15\,ms per turn.
This compares favourably with LLM-as-classifier approaches (50--150\,ms
network round-trip plus API cost) and trained classifiers
(20--40\,ms for a DistilBERT-style forward pass~\cite{routellm},
plus the annotation and retraining pipeline).

\subsection{Conservative escalation policy}

When the maximum cosine similarity across all exemplars falls below
\(\tau\), the router returns opus unconditionally.  The design reflects
the asymmetric cost of misrouting in a coding-agent context: a false
tier-down routes an architectural decision to a weaker model, risking
a broken commit or misarchitected module; a false tier-up routes a
status query to opus, wasting tokens but producing a correct answer.
Given this asymmetry, the router is calibrated to minimise
false-negative rate on the \emph{should-be-opus} class at the cost of
occasional over-spend.

In LOO-CV over the 49-turn evaluation set (§\ref{sec:eval}),
the escalation hatch (\(\tau = 0.30\)) fired on zero turns, meaning
all 49 turns found at least one exemplar with cosine similarity above
0.30.  The hatch guards against truly out-of-distribution queries
(unfamiliar natural-language domains) rather than frequent inputs.

\subsection{Why semantic similarity for v1}

Four considerations favour the exemplar-similarity approach over a
trained classifier at this stage:
(i) \textbf{faster iteration} — adding or correcting a mislabelled
    exemplar requires editing one YAML record, not retraining;
(ii) \textbf{interpretable failures} — a misroute traces to a specific
    exemplar, enabling direct correction;
(iii) \textbf{lower overhead} — one embedder pass plus a dot-product
    is cheaper than any classifier forward pass;
(iv) \textbf{label efficiency} — 50 exemplars per tier is plausible
    for k-NN when classes are geometrically clustered (haiku turns
    cluster around slash-command invocations; sonnet turns cluster
    around implementation requests).
The trained-classifier path (RouteLLM-style~\cite{routellm}) is the
natural v2 direction as the annotation corpus grows (§\ref{sec:future}).

\subsection{Theoretical remark}

Under standard k-NN analysis, if the exemplar set \(E\) is an
i.i.d.\ sample from the operator-turn distribution and class boundaries
are well-separated in embedding space, the expected false-tier-down
rate is bounded by \(1 - \text{Acc}_{k\text{-NN}}\).  In our
LOO-CV evaluation (§\ref{sec:eval}), accuracy is 61.2\,\%, implying
an aggregate misclassification rate of 38.8\,\% — driven primarily by
the opus class (F1 0.125), whose 12 of 49 turns are distributionally
close to sonnet.  Haiku (F1 0.786) and sonnet (F1 0.667) are
better separated.  Closing the opus-recall gap is the primary v2 target
(§\ref{sec:future}).
